\chapter{Structures as categories}
\label{ch:categories}

A structure declared with \texttt{declare-structure}
(Chapter~\ref{ch:defs}) comes with its morphisms: the maps between two
of its models that preserve the structure.  The models and their
morphisms form a category.  This chapter describes how VNB sets those
categories up, how a species can carry more than one, and how one
structure is viewed as another.  Such a view is a functor.

%%% -----------------------------------------------------------------------
\section{Morphisms of structures}
\label{sec:morphisms}

Declaring a structure declares its morphisms too.  From the slot list alone,
\texttt{declare-structure} generates a homomorphism predicate
\texttt{is-hom-\textit{name}}$(a, b, f_1 \ldots f_k)$ --- one map $f_i$ per
\emph{independent} carrier --- whose meaning is the obvious one: each map is
typed between the corresponding carriers, and every operation and constant is
preserved.

\begin{itemize}
\item \textbf{carrier} $c$: a map \texttt{$f$ in fun($c(a)$, $c(b)$)}.
\item \textbf{operation} \texttt{OP}: preservation ---
  \texttt{$f$(OP($a$)($x_1,\ldots$)) = OP($b$)($f(x_1),\ldots$)} for all
  arguments in their sorts (each argument mapped by its own carrier's map).
\item \textbf{constant} \texttt{C}: \texttt{$f$(C($a$)) = C($b$)}.
\item \textbf{substructure} \texttt{K}: agreement, \texttt{K($a$) = K($b$)} ---
  morphisms are taken between structures over the \emph{same} base.
\item \textbf{derived} carrier: nothing.  Its value is pinned by
  \texttt{is-\textit{name}} (Section~\ref{sec:declare-structure}), so it rides
  the base carrier's map and owes no map of its own.  This is why a field
  homomorphism is one map, not two.
\end{itemize}

\noindent A group's generated hom, for instance, is the single map \texttt{$f$
in fun(carr($a$), carr($b$))} with \texttt{$f$(opr($a$)($x$,$y$)) =
opr($b$)($f(x)$,$f(y)$)} and \texttt{$f$(iden($a$)) = iden($b$)} --- exactly a
group homomorphism.  The definition is installed with \emph{definitional}
provenance, so it costs no warrant.

\subsection{Overriding the default: \texttt{declare-hom!}}
\label{sec:declare-hom}

The generated clause for each slot kind is the standard one: preservation for
an operation or a constant, the pullback condition for a \texttt{family} slot
(the arrows of \texttt{top-space} are the continuous maps, those of
\texttt{measurable-space} the measurable maps), transport for a
\texttt{family-fun} slot, ``related points go to related points'' for a
\texttt{relation} slot.  For a \texttt{metric-space} the generated arrows
preserve the distance: they are the isometries, and a \texttt{notation!}
beside the structure says so, since ``homomorphism'' is then the wrong word
(Section~\ref{sec:notation}).  Where no generated clause is the right one, a
species declares its own morphism notion with \texttt{declare-hom!}, beside the
structure like \texttt{notation!}:

\begin{VNBcode}
(declare-hom! 'METRIZABLE-TOP-SPACE '(s t f)
  '(FORALL u (IMPLIES (IN u (OPENS t))
      (IN (PREIMAGE s f u) (OPENS s)))))
\end{VNBcode}

The argument list is the two objects and one map per carrier.  The body states
\emph{only} what is characteristic: the object typings
\texttt{is-metrizable-top-space($s$)}, \texttt{is-metrizable-top-space($t$)}
and the map typing \texttt{$f$ in fun(pts($s$), pts($t$))} are supplied
automatically and need not be repeated.  The declaration replaces the generated
\texttt{is-hom-metrizable-top-space} entirely.  Two species override today:
\texttt{metrizable-top-space} (a refinement inherits its parent's shape, not
its arrows, so continuity is restated) and \texttt{ringoid} (a ring
homomorphism of the underlying rings that preserves the ideal).  The table of
Section~\ref{sec:declare-category} says, for every structure, which of the
three cases its arrows are.

\subsection{Several categories on one species: \texttt{declare-category!}}
\label{sec:declare-category}

A structure is a category: its objects are its models, its arrows the maps
satisfying \texttt{is-hom-\textit{name}}, composed as functions.  That is the
species' \emph{default} category, and for an algebraic species it is the right
one.  For a metric space the default is the category of isometries, which is
not the one analysis works in: the Lipschitz maps, the uniformly continuous
maps and the continuous maps each form a category on the same objects; on the
countably-metrised spaces (\texttt{c-metric-space}) the continuous maps for
the canonical metric form one, which rests on the canonical metric being a
metric (\texttt{c-metric-is-metric-space}).  A species may therefore carry
several categories.  A second one is declared
beside the structure:

\begin{VNBcode}
(declare-category! 'LIPSCHITZ 'METRIC-SPACE '(a b f)
  '(FORSOME lpk_ (AND (IN lpk_ RR) (AND (<= 0 lpk_)
     (FORALL lpx_ (IMPLIES (IN lpx_ (PTS a))
       (FORALL lpy_ (IMPLIES (IN lpy_ (PTS a))
         (<= ((DIST b) (f lpx_) (f lpy_))
             (* lpk_ ((DIST a) lpx_ lpy_)))))))))))
\end{VNBcode}

The body states what is characteristic of an arrow $f \colon a \to b$; the
object typings and the carrier typing of each map are supplied, as for
\texttt{declare-hom!}.  The declaration installs, definitionally, the arrow
predicate \texttt{is-lipschitz-arrow}$(a, b, f)$ and the hom-set
\texttt{hom-lipschitz}$(a, b)$, and registers three \emph{obligations} as
statements, not as theorems: \texttt{hom-lipschitz-id} (the identity is an
arrow), \texttt{hom-lipschitz-compose} (arrows compose) and
\texttt{hom-lipschitz-in-set} (the hom-set is a set).  That the arrows of a
declared category contain the identities and compose is where the
mathematics is (Lipschitz constants multiply; an $\varepsilon$--$\delta$
argument for the continuous maps), so the generator does not stamp it: the
obligations are proven in \texttt{theorem-library/categories.scm}, and
\texttt{category-obligations-audit}, run at the end of every load, lists
every obligation that is not a proven theorem of exactly its statement.
Until that list is empty the declared category is not yet a category.  The
default category is untouched by the declaration.

The categories the library knows are listed below.  The listing is generated
at every load (\texttt{reference/CATEGORIES.md}, also a page of the browser
reference): the default category of each structure with the status of its
three laws, the declared categories with their defining equivalences, and the
inclusions between categories on the same objects --- an isometry is
$1$-Lipschitz, a Lipschitz map is uniformly continuous, a uniformly
continuous map is continuous, a linear isometry is bounded linear --- which
the page finds by the shape of the installed theorems.

%% AUTO-GENERATED by docs/gen-categories.py from reference/CATEGORIES.md -- do not edit.
%% Regenerate:  python3 docs/gen-categories.py   (after a prover load refreshes CATEGORIES.md)

% the long identifiers break only with loose spacing; \sloppy, scoped to this file
\begingroup\sloppy
\subsubsection{The default categories}
One per structure with a morphism predicate.  The laws column reads ``proven'' when the identity law, the composition law and the sethood of the hom-set are all proven theorems; otherwise it names the law and its status.

\begin{longtable}{@{}>{\raggedright\arraybackslash}p{0.22\linewidth}>{\raggedright\arraybackslash}p{0.30\linewidth}>{\raggedright\arraybackslash}p{0.20\linewidth}>{\raggedright\arraybackslash}p{0.18\linewidth}@{}}
\toprule
structure & arrows & defined & laws \\
\midrule
\endhead
\texttt{abelian-\allowbreak{}group} & f is a homomorphism from a to b & generated & proven \\
\texttt{bundle} & f is a homomorphism from a to b & generated & id: not installed; compose: not installed; in-set: not installed \\
\texttt{c-\allowbreak{}metric-\allowbreak{}space} & f is a homomorphism from a to b & generated & proven \\
\texttt{comm-\allowbreak{}monoid} & f is a homomorphism from a to b & generated & proven \\
\texttt{commutative-\allowbreak{}banach-\allowbreak{}algebra} & f is a homomorphism from a to b & refinement of \texttt{unital-\allowbreak{}banach-\allowbreak{}algebra} & id: not installed; compose: not installed; in-set: not installed \\
\texttt{commutative-\allowbreak{}c-\allowbreak{}star-\allowbreak{}algebra} & f is a homomorphism from a to b & refinement of \texttt{unital-\allowbreak{}c-\allowbreak{}star-\allowbreak{}algebra} & id: not installed; compose: not installed; in-set: not installed \\
\texttt{commutative-\allowbreak{}ring} & f is a homomorphism from a to b & refinement of \texttt{ring} & proven \\
\texttt{complex-\allowbreak{}banach-\allowbreak{}algebra} & f is a homomorphism from a to b & refinement of \texttt{complex-\allowbreak{}normed-\allowbreak{}algebra} & id: not installed; compose: not installed; in-set: not installed \\
\texttt{complex-\allowbreak{}banach-\allowbreak{}space} & f is a homomorphism from a to b & refinement of \texttt{complex-\allowbreak{}normed-\allowbreak{}vector-\allowbreak{}space} & id: not installed; compose: not installed; in-set: not installed \\
\texttt{complex-\allowbreak{}inner-\allowbreak{}product-\allowbreak{}space} & f is a homomorphism from a to b & generated & proven \\
\texttt{complex-\allowbreak{}normed-\allowbreak{}algebra} & f is a homomorphism from a to b & generated & id: not installed; compose: not installed; in-set: not installed \\
\texttt{complex-\allowbreak{}normed-\allowbreak{}vector-\allowbreak{}space} & f is a homomorphism from a to b & generated & id: not installed; compose: not installed; in-set: not installed \\
\texttt{complex-\allowbreak{}pre-\allowbreak{}hilbert-\allowbreak{}space} & f is a homomorphism from a to b & refinement of \texttt{complex-\allowbreak{}semi-\allowbreak{}inner-\allowbreak{}product-\allowbreak{}space} & id: not installed; compose: not installed; in-set: not installed \\
\texttt{complex-\allowbreak{}semi-\allowbreak{}inner-\allowbreak{}product-\allowbreak{}space} & f is a homomorphism from a to b & generated & id: not installed; compose: not installed; in-set: not installed \\
\texttt{directed-\allowbreak{}set} & f is a homomorphism from a to b & generated & id: not installed; compose: not installed; in-set: not installed \\
\texttt{euclidean-\allowbreak{}ring} & f is a homomorphism from a to b & refinement of \texttt{integral-\allowbreak{}domain} & proven \\
\texttt{field} & f is a homomorphism from a to b & generated & proven \\
\texttt{field-\allowbreak{}ring} & f is a homomorphism from a to b & refinement of \texttt{commutative-\allowbreak{}ring} & proven \\
\texttt{group} & f is a homomorphism from a to b & generated & proven \\
\texttt{infiltrated-\allowbreak{}space} & f1 is a morphism of infiltrated strand spaces from a to b & generated & id: not installed; compose: not installed; in-set: not installed \\
\texttt{integral-\allowbreak{}domain} & f is a homomorphism from a to b & refinement of \texttt{commutative-\allowbreak{}ring} & proven \\
\texttt{measurable-\allowbreak{}space} & f is a measurable map from a to b & generated & proven \\
\texttt{measure-\allowbreak{}space} & f is a measure-preserving measurable map from a to b & generated & proven \\
\texttt{metric-\allowbreak{}group} & f is an isometric homomorphism from a to b & generated & id: not installed; compose: not installed; in-set: not installed \\
\texttt{metric-\allowbreak{}space} & f is an isometry from a to b & generated & proven \\
\texttt{metrizable-\allowbreak{}top-\allowbreak{}space} & f is continuous from a to b & by \texttt{declare-\allowbreak{}hom!} & proven \\
\texttt{module} & f is a homomorphism from a to b & generated & proven \\
\texttt{monoid} & f is a homomorphism from a to b & generated & proven \\
\texttt{normed-\allowbreak{}ag} & f is a homomorphism from a to b & generated & proven \\
\texttt{normed-\allowbreak{}field} & f is a homomorphism from a to b & generated & proven \\
\texttt{normed-\allowbreak{}vector-\allowbreak{}space} & f is a homomorphism from a to b & generated & proven \\
\texttt{nu-\allowbreak{}c-\allowbreak{}star-\allowbreak{}algebra} & f is a homomorphism from a to b & refinement of \texttt{nu-\allowbreak{}involutive-\allowbreak{}banach-\allowbreak{}algebra} & id: not installed; compose: not installed; in-set: not installed \\
\texttt{nu-\allowbreak{}involutive-\allowbreak{}banach-\allowbreak{}algebra} & f is a homomorphism from a to b & generated & id: not installed; compose: not installed; in-set: not installed \\
\texttt{pid} & f is a homomorphism from a to b & refinement of \texttt{integral-\allowbreak{}domain} & proven \\
\texttt{pk-\allowbreak{}space} & f is a morphism of public-key strand spaces from a to b & generated & id: not installed; compose: not installed; in-set: not installed \\
\texttt{pre-\allowbreak{}hilbert-\allowbreak{}space} & f is a homomorphism from a to b & refinement of \texttt{semi-\allowbreak{}inner-\allowbreak{}product-\allowbreak{}space} & id: not installed; compose: not installed; in-set: not installed \\
\texttt{pseudometric-\allowbreak{}space} & f is a homomorphism from a to b & generated & proven \\
\texttt{ring} & f is a homomorphism from a to b & generated & proven \\
\texttt{ringoid} & f is a ringoid homomorphism from a to b & by \texttt{declare-\allowbreak{}hom!} & proven \\
\texttt{semi-\allowbreak{}inner-\allowbreak{}product-\allowbreak{}space} & f is a homomorphism from a to b & generated & id: not installed; compose: not installed; in-set: not installed \\
\texttt{semigroup} & f is a homomorphism from a to b & generated & proven \\
\texttt{setoid} & f is a setoid morphism from a to b & generated & proven \\
\texttt{sk-\allowbreak{}space} & f is a morphism of shared-key strand spaces from a to b & generated & id: not installed; compose: not installed; in-set: not installed \\
\texttt{strand-\allowbreak{}space} & f1 is a morphism of strand spaces from a to b & generated & id: not installed; compose: not installed; in-set: not installed \\
\texttt{top-\allowbreak{}group} & f is a continuous homomorphism from a to b & generated & id: not installed; compose: not installed; in-set: not installed \\
\texttt{top-\allowbreak{}space} & f is continuous from a to b & generated & proven \\
\texttt{unital-\allowbreak{}banach-\allowbreak{}algebra} & f is a homomorphism from a to b & generated & id: not installed; compose: not installed; in-set: not installed \\
\texttt{unital-\allowbreak{}c-\allowbreak{}star-\allowbreak{}algebra} & f is a homomorphism from a to b & refinement of \texttt{unital-\allowbreak{}involutive-\allowbreak{}banach-\allowbreak{}algebra} & id: not installed; compose: not installed; in-set: not installed \\
\texttt{unital-\allowbreak{}involutive-\allowbreak{}banach-\allowbreak{}algebra} & f is a homomorphism from a to b & generated & id: not installed; compose: not installed; in-set: not installed \\
\texttt{vector-\allowbreak{}space} & f is a homomorphism from a to b & refinement of \texttt{module} & proven \\
\bottomrule
\end{longtable}

\subsubsection{The declared categories}

\paragraph{\texttt{bounded-\allowbreak{}linear,\allowbreak{} on normed-\allowbreak{}vector-\allowbreak{}space}}
\emph{Declared in} \texttt{theorem-\allowbreak{}library/\allowbreak{}categories.scm}.

Objects: the models of \texttt{is-\allowbreak{}normed-\allowbreak{}vector-\allowbreak{}space}.  Arrows: \texttt{is-\allowbreak{}bounded-\allowbreak{}linear-\allowbreak{}arrow(a,\allowbreak{} b,\allowbreak{} f)} --- f is a bounded linear map from a to b.  Hom-set: \texttt{hom-\allowbreak{}bounded-\allowbreak{}linear(a,\allowbreak{} b)} --- the bounded linear maps from a to b.

Definition (\texttt{is-\allowbreak{}bounded-\allowbreak{}linear-\allowbreak{}arrow-\allowbreak{}def}, definitional):
\begin{VNBsnip}
forall([a, b, f],
  is-bounded-linear-arrow(a, b, f)
  iff is-normed-vector-space(a)
      and is-normed-vector-space(b)
      and f in fun(vec(a), vec(b))
      and scal(a) = scal(b)
      and forall([x1_ in vec(a), x2_ in vec(a)],
            f((vadd(a))(x1_, x2_)) = (vadd(b))(f(x1_), f(x2_)))
      and forall([x1_ in carr(scal(a)), x2_ in vec(a)],
            f((act(a))(x1_, x2_)) = (act(b))(x1_, f(x2_)))
      and forsome([blk_ in rr],
            0 <= blk_
            and forall([x1_ in vec(a)],
                  (vnrm(b))(f(x1_)) <= blk_ * (vnrm(a))(x1_))))
\end{VNBsnip}

Laws: \texttt{hom-\allowbreak{}bounded-\allowbreak{}linear-\allowbreak{}id} --- proven; \texttt{hom-\allowbreak{}bounded-\allowbreak{}linear-\allowbreak{}compose} --- proven; \texttt{hom-\allowbreak{}bounded-\allowbreak{}linear-\allowbreak{}in-\allowbreak{}set} --- proven; membership \texttt{hom-\allowbreak{}bounded-\allowbreak{}linear-\allowbreak{}member-\allowbreak{}iff} --- proven; Hom functors \texttt{hom-\allowbreak{}bounded-\allowbreak{}linear-\allowbreak{}post-\allowbreak{}type} --- proven, \texttt{hom-\allowbreak{}bounded-\allowbreak{}linear-\allowbreak{}pre-\allowbreak{}type} --- proven.

\paragraph{\texttt{c-\allowbreak{}continuous,\allowbreak{} on c-\allowbreak{}metric-\allowbreak{}space}}
\emph{Declared in} \texttt{theorem-\allowbreak{}library/\allowbreak{}c-\allowbreak{}continuous-\allowbreak{}category.scm}.

Objects: the models of \texttt{is-\allowbreak{}c-\allowbreak{}metric-\allowbreak{}space}.  Arrows: \texttt{is-\allowbreak{}c-\allowbreak{}continuous-\allowbreak{}arrow(a,\allowbreak{} b,\allowbreak{} f)} --- f is a continuous map of countably-metrised spaces from a to b.  Hom-set: \texttt{hom-\allowbreak{}c-\allowbreak{}continuous(a,\allowbreak{} b)} --- the continuous maps of countably-metrised spaces from a to b.

Definition (\texttt{is-\allowbreak{}c-\allowbreak{}continuous-\allowbreak{}arrow-\allowbreak{}def}, definitional):
\begin{VNBsnip}
forall([a, b, f],
  is-c-continuous-arrow(a, b, f)
  iff is-c-metric-space(a)
      and is-c-metric-space(b)
      and f in fun(pts(a), pts(b))
      and is-continuous(c-metric(a), c-metric(b), f))
\end{VNBsnip}

Laws: \texttt{hom-\allowbreak{}c-\allowbreak{}continuous-\allowbreak{}id} --- proven; \texttt{hom-\allowbreak{}c-\allowbreak{}continuous-\allowbreak{}compose} --- proven; \texttt{hom-\allowbreak{}c-\allowbreak{}continuous-\allowbreak{}in-\allowbreak{}set} --- proven; membership \texttt{hom-\allowbreak{}c-\allowbreak{}continuous-\allowbreak{}member-\allowbreak{}iff} --- proven; Hom functors \texttt{hom-\allowbreak{}c-\allowbreak{}continuous-\allowbreak{}post-\allowbreak{}type} --- proven, \texttt{hom-\allowbreak{}c-\allowbreak{}continuous-\allowbreak{}pre-\allowbreak{}type} --- proven.

\paragraph{\texttt{continuous,\allowbreak{} on metric-\allowbreak{}space}}
\emph{Declared in} \texttt{theorem-\allowbreak{}library/\allowbreak{}categories.scm}.

Objects: the models of \texttt{is-\allowbreak{}metric-\allowbreak{}space}.  Arrows: \texttt{is-\allowbreak{}continuous-\allowbreak{}arrow(a,\allowbreak{} b,\allowbreak{} f)} --- f is a continuous map from a to b.  Hom-set: \texttt{hom-\allowbreak{}continuous(a,\allowbreak{} b)} --- the continuous maps from a to b.

Definition (\texttt{is-\allowbreak{}continuous-\allowbreak{}arrow-\allowbreak{}def}, definitional):
\begin{VNBsnip}
forall([a, b, f],
  is-continuous-arrow(a, b, f)
  iff is-metric-space(a)
      and is-metric-space(b)
      and f in fun(pts(a), pts(b))
      and is-continuous(a, b, f))
\end{VNBsnip}

Laws: \texttt{hom-\allowbreak{}continuous-\allowbreak{}id} --- proven; \texttt{hom-\allowbreak{}continuous-\allowbreak{}compose} --- proven; \texttt{hom-\allowbreak{}continuous-\allowbreak{}in-\allowbreak{}set} --- proven; membership \texttt{hom-\allowbreak{}continuous-\allowbreak{}member-\allowbreak{}iff} --- proven; Hom functors \texttt{hom-\allowbreak{}continuous-\allowbreak{}post-\allowbreak{}type} --- proven, \texttt{hom-\allowbreak{}continuous-\allowbreak{}pre-\allowbreak{}type} --- proven.

\paragraph{\texttt{lipschitz,\allowbreak{} on metric-\allowbreak{}space}}
\emph{Declared in} \texttt{theorem-\allowbreak{}library/\allowbreak{}categories.scm}.

Objects: the models of \texttt{is-\allowbreak{}metric-\allowbreak{}space}.  Arrows: \texttt{is-\allowbreak{}lipschitz-\allowbreak{}arrow(a,\allowbreak{} b,\allowbreak{} f)} --- f is a Lipschitz map from a to b.  Hom-set: \texttt{hom-\allowbreak{}lipschitz(a,\allowbreak{} b)} --- the Lipschitz maps from a to b.

Definition (\texttt{is-\allowbreak{}lipschitz-\allowbreak{}arrow-\allowbreak{}def}, definitional):
\begin{VNBsnip}
forall([a, b, f],
  is-lipschitz-arrow(a, b, f)
  iff is-metric-space(a)
      and is-metric-space(b)
      and f in fun(pts(a), pts(b))
      and forsome([lpk_ in rr],
            0 <= lpk_
            and forall([lpx_ in pts(a), lpy_ in pts(a)],
                  (dist(b))(f(lpx_), f(lpy_)) <= lpk_ * (dist(a))(lpx_, lpy_))))
\end{VNBsnip}

Laws: \texttt{hom-\allowbreak{}lipschitz-\allowbreak{}id} --- proven; \texttt{hom-\allowbreak{}lipschitz-\allowbreak{}compose} --- proven; \texttt{hom-\allowbreak{}lipschitz-\allowbreak{}in-\allowbreak{}set} --- proven; membership \texttt{hom-\allowbreak{}lipschitz-\allowbreak{}member-\allowbreak{}iff} --- proven; Hom functors \texttt{hom-\allowbreak{}lipschitz-\allowbreak{}post-\allowbreak{}type} --- proven, \texttt{hom-\allowbreak{}lipschitz-\allowbreak{}pre-\allowbreak{}type} --- proven.

\paragraph{\texttt{uniformly-\allowbreak{}continuous,\allowbreak{} on metric-\allowbreak{}space}}
\emph{Declared in} \texttt{theorem-\allowbreak{}library/\allowbreak{}categories.scm}.

Objects: the models of \texttt{is-\allowbreak{}metric-\allowbreak{}space}.  Arrows: \texttt{is-\allowbreak{}uniformly-\allowbreak{}continuous-\allowbreak{}arrow(a,\allowbreak{} b,\allowbreak{} f)} --- f is a uniformly continuous map from a to b.  Hom-set: \texttt{hom-\allowbreak{}uniformly-\allowbreak{}continuous(a,\allowbreak{} b)} --- the uniformly continuous maps from a to b.

Definition (\texttt{is-\allowbreak{}uniformly-\allowbreak{}continuous-\allowbreak{}arrow-\allowbreak{}def}, definitional):
\begin{VNBsnip}
forall([a, b, f],
  is-uniformly-continuous-arrow(a, b, f)
  iff is-metric-space(a)
      and is-metric-space(b)
      and f in fun(pts(a), pts(b))
      and is-uniformly-continuous(a, b, f))
\end{VNBsnip}

Laws: \texttt{hom-\allowbreak{}uniformly-\allowbreak{}continuous-\allowbreak{}id} --- proven; \texttt{hom-\allowbreak{}uniformly-\allowbreak{}continuous-\allowbreak{}compose} --- proven; \texttt{hom-\allowbreak{}uniformly-\allowbreak{}continuous-\allowbreak{}in-\allowbreak{}set} --- proven; membership \texttt{hom-\allowbreak{}uniformly-\allowbreak{}continuous-\allowbreak{}member-\allowbreak{}iff} --- proven; Hom functors \texttt{hom-\allowbreak{}uniformly-\allowbreak{}continuous-\allowbreak{}post-\allowbreak{}type} --- proven, \texttt{hom-\allowbreak{}uniformly-\allowbreak{}continuous-\allowbreak{}pre-\allowbreak{}type} --- proven.

\subsubsection{Inclusions}
Theorems of the form $P(a, b, f) \Rightarrow Q(a, b, f)$ for two arrow predicates, or $H_1(a, b) \subseteq H_2(a, b)$ for two hom-sets, found by shape.

\begin{itemize}
\item \texttt{hom-\allowbreak{}metric-\allowbreak{}space-\allowbreak{}is-\allowbreak{}lipschitz-\allowbreak{}arrow} (proven): \texttt{is-\allowbreak{}hom-\allowbreak{}metric-\allowbreak{}space(a,\allowbreak{} b,\allowbreak{} f) implies is-\allowbreak{}lipschitz-\allowbreak{}arrow(a,\allowbreak{} b,\allowbreak{} f)}
\item \texttt{hom-\allowbreak{}normed-\allowbreak{}vector-\allowbreak{}space-\allowbreak{}is-\allowbreak{}bounded-\allowbreak{}linear-\allowbreak{}arrow} (proven): \texttt{is-\allowbreak{}hom-\allowbreak{}normed-\allowbreak{}vector-\allowbreak{}space(a,\allowbreak{} b,\allowbreak{} f) implies is-\allowbreak{}bounded-\allowbreak{}linear-\allowbreak{}arrow(a,\allowbreak{} b,\allowbreak{} f)}
\item \texttt{lipschitz-\allowbreak{}arrow-\allowbreak{}is-\allowbreak{}uniformly-\allowbreak{}continuous-\allowbreak{}arrow} (proven): \texttt{is-\allowbreak{}lipschitz-\allowbreak{}arrow(a,\allowbreak{} b,\allowbreak{} f) implies is-\allowbreak{}uniformly-\allowbreak{}continuous-\allowbreak{}arrow(a,\allowbreak{} b,\allowbreak{} f)}
\item \texttt{uniformly-\allowbreak{}continuous-\allowbreak{}arrow-\allowbreak{}is-\allowbreak{}continuous-\allowbreak{}arrow} (proven): \texttt{is-\allowbreak{}uniformly-\allowbreak{}continuous-\allowbreak{}arrow(a,\allowbreak{} b,\allowbreak{} f) implies is-\allowbreak{}continuous-\allowbreak{}arrow(a,\allowbreak{} b,\allowbreak{} f)}
\end{itemize}
\par\endgroup


\section{Viewing one structure as another: \texttt{view-as}}
\label{sec:view-as}

A ring's additive part $(\mathtt{carr}, \mathtt{add}, \mathtt{zero},
\mathtt{neg})$ is an abelian group; its multiplicative part
$(\mathtt{carr}, \mathtt{mul}, \mathtt{one})$ is a monoid.  A
\emph{view-as} records such a relationship: a way of \emph{seeing} an
instance of one structure as an instance of another.  The relationship
is asymmetric (the ring's additive group forgets the multiplication;
there is no route back) and not unique (the same ring has both an
additive-group view and a multiplicative-monoid view).  We call these
\emph{views}, not ``bridges'': a bridge sounds symmetric, and the
asymmetry and the multiplicity are exactly the point.

A view is a \emph{functoid} (Remark~\ref{rem:functoids-quantified}):
\texttt{def-functor} installs the \texttt{def-functoid} that builds the
target instance as the list of the source's slots, plus the definitional
axiom that the result satisfies the target predicate.  Its domain, the
class of all instances of the source structure, is a proper class; a
view is a member of nothing and, like every functoid, outside the range
of the quantifiers.

\subsection{The \texttt{def-functor} form}

\begin{VNBcode}
(def-functor NAME
  SOURCE-STRUCT  (src-comp-1 src-comp-2 ... src-comp-n)
  TARGET-STRUCT  (tgt-slot-1 tgt-slot-2 ... tgt-slot-n))
\end{VNBcode}

The two component lists are flat lists of accessor symbols of equal
length, matched position by position: the value at the target's $k$-th
slot is the source's $k$-th named component.  The target list must be
the target structure's slots in their declared order (the form checks
this); the source list names, in that same order, which source accessor
fills each target slot.  So a view is a \textbf{slot map written in the
target's coordinate system}.

A view-as is a \emph{forgetful functor} between structure categories, and a
forgetful functor does exactly one or both of two things (recall
\S\ref{sec:shape-vs-property}: a structure is a \emph{shape} plus
\emph{properties}):

\begin{itemize}
\item \textbf{forget shape} --- the source list keeps only a sub-list of the
  source's slots; the rest drop.  (A ring's additive group forgets
  \texttt{mul}/\texttt{one}; a field's multiplicative group keeps the
  \texttt{non-zero} carrier slot and forgets the additive ones.)
\item \textbf{forget properties} --- \texttt{is-\textit{target}} folds fewer laws than
  \texttt{is-\textit{source}}, so the image satisfies less (a commutative ring viewed
  as a ring forgets commutativity).
\end{itemize}

\noindent These are the \emph{only} two forgetful operations.  Writing the kept
slots in the target's coordinate order (a ring's \texttt{add} filling a group's
\texttt{mul} slot, or a non-first slot serving as the carrier) is just
notation, not a third kind --- it is what lets shape-forgetting be expressed at
all.  Genuinely \emph{non}-forgetful functors --- the quotient, the completion
--- are constructions, not views; they are written as plain functoids.

\begin{VNBcode}
;;; A ring's additive structure is an abelian group.
;;; ABELIAN-GROUP's slots, in order, are (CARR OPR IDEN INV); the view fills
;;; them from the ring's (CARR ADD ZERO NEG):
(def-functor 'RING-ADDITIVE-AG
  'RING          '(CARR ADD ZERO NEG)
  'ABELIAN-GROUP '(CARR OPR IDEN INV))
\end{VNBcode}

\subsection{What \texttt{def-functor} installs}

\begin{enumerate}
\item \textbf{A constructor functoid} \texttt{(\textit{name} s)} whose defining
  equation gathers the named source components into a target-shaped list:
  \[
    \mathtt{(ring\text{-}additive\text{-}ag}\;r) \;=\;
      \mathtt{list}(\mathtt{carr}(r),\, \mathtt{add}(r),\,
                    \mathtt{zero}(r),\, \mathtt{neg}(r)).
  \]
  Inside a proof, \texttt{(mul (ring-additive-ag r))} reduces to
  \texttt{(add r)} by unfolding the functoid and the accessor.

\item \textbf{A typing axiom}, named
  \texttt{\textit{name}-is-\textit{target}} (here
  \texttt{ring-additive-ag-is-abelian-group}):
  \[
    \forall r.\; \mathit{is\text{-}source}(r) \;\Rightarrow\;
      \mathit{is\text{-}target}(\mathit{name}(r)).
  \]
  During the library-build phase this is installed as an \emph{axiom}
  (a warranted assertion, Section~\ref{sec:warrants}), not proved;
  promoting it to a theorem is a separate task.

\item \textbf{Auto-specializations.}  Every theorem of the form
  $\forall s.\; \mathit{is\text{-}target}(s) \Rightarrow P[s]$ already in
  the library is re-installed as
  $\forall r.\; \mathit{is\text{-}source}(r) \Rightarrow P[\mathit{name}(r)]$,
  with the target-slot accessors reduced to source components (so
  $\mathtt{mul}(s)$ becomes $\mathtt{add}(r)$, etc.), under the name
  \texttt{\textit{original}-\textit{name}}.  Thus declaring the additive
  view immediately carries every abelian-group theorem over to the
  additive structure of every ring.
\end{enumerate}

\begin{remark}
A view-as is a covariant functor in waiting: the constructor is the
object map (an instance of the source becomes an instance of the
target), and the auto-specialization transports the target's theorems
along it.  Homomorphisms are now first-class
(Section~\ref{sec:morphisms}); the morphism half --- transporting a
homomorphism across a view --- remains deferred.  The
auto-specializer currently matches only theorems guarded by exactly
\texttt{is-\textit{target}}; it does not chain through structure-inclusion axioms,
so same-shape refinements (a commutative ring's additive view, say) are
declared directly rather than derived.
\end{remark}

\begin{remark}
Same-shape inclusions --- every abelian group is a group, every
commutative ring is a ring --- are \emph{not} views.  They are recorded
as plain $\texttt{is-}X \Rightarrow \texttt{is-}Y$ typing axioms, because
the constructor would be a redundant identity list.  A \texttt{view-as}
earns its keep when source and target have different shapes, so the
constructor genuinely projects.
\end{remark}

\subsection{Constructed functors: \texttt{def-constructed-functor}}
\label{sec:constructed-functor}

A \texttt{view-as} selects the target's slots from among the source's --- its
object map is a re-shuffling of accessors.  Some functors are not like that.
The metric topology $\tau_d$ is \emph{computed} from the distance, not read off a
slot; the distance of a normed field is \emph{computed} from the norm.  For these
the object map is an arbitrary term, and the form is
\texttt{def-constructed-functor}:

\begin{VNBcode}
;;; A normed field is a metric space, with d(x, y) = norm(x - y).
(def-constructed-functor 'NF-METRIC-SPACE 'NORMED-FIELD 'METRIC-SPACE '(nf)
  '(LIST (CARR nf)
         (VNB-LAMBDA (LIST x y) (CARTESIAN (CARR nf) (CARR nf)) ((FNRM nf) ((ADD nf) x ((NEG nf) y))))))
\end{VNBcode}

The arguments are the functor name, its source and target species, the
parameter list, and the defining term --- installed as an ordinary functoid
(Section~\ref{sec:def-functoid}), so \texttt{(nf-metric-space nf)} rewrites to
that tuple.

The decisive difference from \texttt{view-as} is what it \emph{claims}.  A view's
typing axiom is installed as a warranted assertion; a constructed functor
\textbf{asserts nothing}.  ``An isometry is continuous'' and ``the norm-difference
is a metric'' are theorems with content, not free consequences of a slot
correspondence.  So the form only \emph{records} the two obligations any functor
owes, and \texttt{functor-obligation-audit} reports at load any still
undischarged:
\[
  \begin{aligned}
    \texttt{\textit{name}-is-\textit{tgt}} &:\quad
      \mathit{is\text{-}src}(r) \Rightarrow \mathit{is\text{-}tgt}(\mathit{name}(r)), \\
    \texttt{\textit{name}-functorial} &:\quad
      \mathit{is\text{-}hom\text{-}src}(a, b, f) \Rightarrow
      \mathit{is\text{-}hom\text{-}tgt}(\mathit{name}(a), \mathit{name}(b), f).
  \end{aligned}
\]
Both are available as goals, ready for \texttt{sp}.  A functor whose obligations
are unproved is a functor you do not yet have.
